\begin{proposition}[an unramified torsor alternative]\hypertarget{tag-030A}{}\label{push:src:B.3.6}\label{push:B36}
\StatementRecord{push:src:B.3.6}{proved}{030A}{\SeriesRef{push:src:B.3.4}, \SeriesRef{push:src:B.4.1}; \cite[Chapter III, Section 4, Theorem 2
and the following remark, pp.129--130]{HS53}; \cite[Lemma 4.7.8]{CDL24}; \cite{CDL24}; \cite[Section 4, equation (4)]{Xiao91}; \cite[Sections 3--4]{Xiao91}}

Put $K=\mathbb C(t)$, $L=\mathbb C(C)$,
$G=\operatorname{Gal}(L/K)=\langle\sigma\rangle$,
and let $E/K$ be the generic elliptic
curve of $S\to B$. There is a nonzero
class
\[
 \xi\in\operatorname{Sha}(E_L/L)^G
 \quad\hbox{whose transgression in}\quad
 H^2(G,E(L))=\mathbb ZP/12\mathbb ZP
 \quad\hbox{is }6P.
\]
Here $\operatorname{Sha}$ denotes the
subgroup of genus-one torsors trivial
over the henselian local fields of all
points of $C$.
In particular the trace obstruction of
\ref{push:src:B.3.4} does not obstruct all genus-one
torsors. Such a class has a relatively
minimal projective genus-one surface
$V\to C$, locally isomorphic to
$\mathscr E\to C$, and an affine lift
$A:V\to V$ of $\sigma$ with
\[
 A^{12}=t_{6P}.
\]
This is an existence statement for a
possible normalization base, not a
construction of a threefold central
fibre or of its smoothing.
\end{proposition}


\begin{proof}
\begin{enumerate}[label=\arabic*.,leftmargin=*]
\item We use Galois cohomology with its
continuous cochains. The five-term
Hochschild--Serre sequence gives
\[
 H^1(K,E)\longrightarrow H^1(L,E)^G
 \xrightarrow{\ d_2\ }H^2(G,E(L))
 \longrightarrow H^2(K,E).
\]
The transgression sequence is
\cite[Chapter III, Section 4, Theorem 2
and the following remark, pp.129--130]{HS53};
its continuous form follows by passing
to finite Galois quotients for the
finitely many values of each continuous
cochain. For this complex function field,
$H^2(K,E)=0$, as in
\cite[Lemma 4.7.8]{CDL24}.
One can also see this from
$\operatorname{cd}K=1$: the torsion
submodule of $E(\overline K)$ has no
second cohomology, and its quotient is
a rational vector space, whose positive
continuous cohomology vanishes by
averaging on finite quotients.

\item Cyclic group cohomology gives
$H^2(G,E(L))=E(K)/\operatorname{Tr}_\sigma E(L)$.
By \ref{push:src:B.3.4} this is $\mathbb Z/12$, with
generator the class of $P$. Hence there
is a $G$-invariant torsor class
$\xi_0\in H^1(L,E)$ with $d_2\xi_0=6P$.
We next remove all its local ramification
without changing this transgression.

\item At a place $v$ of $B$ choose a point
$w$ of $C$ above it and its decomposition
subgroup $G_w$. The only nontrivial
decomposition groups have orders
$3,4,12$, at $0,1,\infty$ respectively.
The component groups of the original
fibres are $\mathbb Z/3$, $\mathbb Z/2$
and zero. Thus $6P$ belongs to the
identity component of the local N\'eron
group at every ramified place.
Multiplication by the ramification
index is surjective on that identity
group over the henselian local field:
on the residue group it is surjective
on $\mathbb G_a(\mathbb C)$ or
$\mathbb G_m(\mathbb C)$, and its
derivative is invertible; Hensel's
lemma lifts a division point.
Choose such a division point over
$K_v^h$. Viewed over $L_w^h$, it is
invariant, so its local norm is $6P$.
The restriction of the transgression
class to $H^2(G_w,E(L_w^h))$ is
therefore zero. For unramified places
the decomposition group is trivial,
so the same assertion is automatic.

\item Naturality and exactness of the local
five-term sequences now show that
$(\xi_0)_w$ is the restriction of some
$\eta_v\in H^1(K_v^h,E)$.
Only finitely many of these classes
need be nonzero. The original elliptic
surface is nontrivial, since
$j=1728t$. Ogg--Shafarevich localization
is therefore surjective:
\[
 H^1(K,E)\longrightarrow
 \bigoplus_{v\in B}H^1(K_v^h,E)
 \longrightarrow0.
\]
This is the combination of exact
sequence (4.7.4) and Theorem 4.7.9
of \cite{CDL24}; the same surjectivity
and its geometric fibre-twist
interpretation are given in
\cite[Section 4, equation (4)]{Xiao91}.
Take a global $\eta$ with these local
values and put
$\xi=\xi_0-\operatorname{res}_{L/K}\eta$.
It is locally trivial at every point
of $C$, remains $G$-invariant, and has
$d_2\xi=6P$. It is nonzero because its
transgression is nonzero.

\item The standard interpretation of the
Weil--Ch\^atelet group represents
$\xi$ by a smooth projective genus-one
curve over $L$; taking its relatively
minimal regular model over $C$ gives
$V$. Local triviality identifies each
henselian model with that of $E_L$.
Over $\mathbb C$ the corresponding
\'etale identifications yield
holomorphic identifications over small
base discs. Equivalently, $V$ is obtained
by gluing copies of $\mathscr E$ over
these discs by translations by local
sections. These translations extend
over the regular $I_{12}$ model.
Thus $V$ is locally isomorphic to
$\mathscr E$, has smooth total surface,
no multiple fibres, and the same single
$I_{12}$ fibre. Properness and flatness
hold locally over the base and hence
globally. This use of torsors and local
twists is the one in
\cite[Sections 3--4]{Xiao91}; it does
not assert an explicit equation for
the selected torsor.

\item The $G$-invariance of $\xi$ gives an
affine isomorphism of its generic fibre
covering $\sigma$. It extends to the
relatively minimal regular surface:
in this particular situation the
extension can be checked on a common
resolution. If an exceptional
$(-1)$-curve for one projection is not
contracted by the other, its image is
a vertical rational curve and has
self-intersection at least $-1$,
since blowing down only increases
self-intersection. On either locally
trivial model all vertical rational
curves are components of the $I_{12}$
fibre, of square $-2$. Thus every such
$(-1)$-curve is contracted by both
projections and can be removed from
the common resolution. Repeating this
proves that the map is an isomorphism.
The same argument gives the local
comparison with the Jacobian model
used above. Transporting the base by
$\sigma$ changes none of these facts.
Write $A_0$ for this lift. Its twelfth
power is translation by a section
$N\in E(L)^G=\mathbb ZP$.
Choosing the powers $1,A_0,\ldots,A_0^{11}$
as lifts of $G$ gives the cyclic
factor set
$N\lfloor(i+j)/12\rfloor$.
Its cohomology class is precisely the
descent obstruction $d_2\xi$, so
$N=(6+12m)P$ for an integer $m$.
Replacing $A_0$ by $A=t_{-mP}\circ A_0$
gives $A^{12}=t_{6P}$.
The finite carry-cocycle identities,
its exact order two modulo twelve,
and the three local component checks
are recorded in
\texttt{ab\_push\_torsor\_cocycles.py}.

\item Neither this cohomological construction
nor local triviality specifies the
normal line bundles, the threefold
cusp charts, or a smoothing with the
given punctured-family identification.
Those remain requirements of \ref{push:src:B.4.1}.
\end{enumerate}
\end{proof}
